\begin{table}[!h]
\centering\centering\centering
\caption{Effect on Absolute Deviation in Test Scores from Reading Group Mean}
\centering
\begin{threeparttable}
\fontsize{10}{12}\selectfont
\begin{tabular}[t]{lccccccc}
\toprule
\multicolumn{1}{c}{ } & \multicolumn{7}{c}{Experiment (Liberia)} \\
\cmidrule(l{3pt}r{3pt}){2-8}
  & (1) & (2) & (3) & (4) & (5) & (6) & (7)\\
\midrule
Peer Predicted Score & -0.030 &  &  &  & -0.026 &  & -0.014\\
 & (0.065) &  &  &  & (0.072) &  & (0.076)\\
Dispersion in Peer Scores &  & 0.533*** &  &  &  & 0.472** & 0.462**\\
 &  & (0.180) &  &  &  & (0.198) & (0.199)\\
Distance from Median Instruction &  &  & -0.039 &  & -0.024 & -0.014 & -0.012\\
 &  &  & (0.065) &  & (0.064) & (0.058) & (0.060)\\
Class Size &  &  &  & 0.004** & 0.005** & 0.003 & 0.003\\
 &  &  &  & (0.002) & (0.002) & (0.002) & (0.002)\\
\hline
\midrule
Number of Students & 2783 & 2783 & 2409 & 3135 & 2408 & 2408 & 2408\\
Number of Schools & 54 & 54 & 54 & 54 & 54 & 54 & 54\\
Number of Test Scores & 1 & 1 & 1 & 1 & 1 & 1 & 1\\
\bottomrule
\end{tabular}
\begin{tablenotes}[para]
\item Notes: All specifications control for individual baseline test score,
                      grade dummies and randomization strata fixed effects. Each specification
                      also controls for the expected value of the respective independent variable
                      used in the specification. Standard errors are clustered at the academy-grade group                        level. ***, **, and * indicate significance at 1\%, 5\%, and 10\%.
\end{tablenotes}
\end{threeparttable}
\end{table}
